\documentclass[output=paper,colorlinks,citecolor=brown]{langscibook}
\ChapterDOI{10.5281/zenodo.20541860}
\author{Robin Meyer\orcid{0000-0002-1177-0639}\affiliation{Université de Lausanne}}
\title[Hypoanalysis and aspectual diversification]{Hypoanalysis and aspectual diversification: The Armenian future in diachrony}
\abstract{
	Old Armenian does not exhibit a separate future tense; future events are commonly
	expressed by means of modal forms. 	In Middle Armenian, the historic subjunctive is lost as part of the process of replacing the synthetic present with an analytical, originally progressive form; the historic present indicative takes on subjunctive functions. Later in this period, a future develops by hypoanalysis and, secondarily, univerbation of
	the former present indicative, whilst an innovative periphrastic form supersedes it; the
	subjunctive persists.\\
	In Modern Eastern Armenian, the new future takes on additional, secondary functions; other, analytical future periphrases have arisen, increasing the aspectual
	differentiation of the language in analogy to other tenses; not all are equally productive. The synthetic future might be considered an aspectually unmarked form.\\
	This paper outlines in as much detail as currently possible the development of the future from Old to Modern Armenian, with particular attention to grammaticalisation and hypoanalysis processes. Furthermore, a corpus-based sketch of the future periphrases in Modern Eastern Armenian is given with a view to proposing a basic diachrony of and motivation for their development.}

\IfFileExists{../localcommands.tex}{
	\addbibresource{localbibliography.bib}
	\usepackage{tabularx,multicol}
\usepackage{url}
\urlstyle{same}

\usepackage[english]{babel}
 
\usepackage{langsci-optional}
\usepackage{langsci-lgr}
\usepackage{langsci-gb4e}

\usepackage{langsci-branding}

%\usepackage{tabularx}

% elena
\usepackage{graphicx} %package to manage images
\usepackage{paralist}

% sarah & lena
\usepackage[justification=centering]{caption}
\usepackage{rotating}

% robin
\usepackage{lscape}
\usepackage{tikz}
\usetikzlibrary{tikzmark}

% dominika
\usepackage{multirow}
\usepackage{hanging}

% olaf & dylan
\usepackage{subcaption}

% svetlana
\usepackage{verbatim}
\usepackage{enumitem}

% wiemer
%\usepackage[markup=underlined]{changes}
%\usepackage{hyphenat}

% wir
\usepackage{longtable}

% for crossref in the introduction jiabin
\usepackage{xr}
	% \newcommand*{\orcid}{}
%

\makeatletter
\let\thetitle\@title
\let\theauthor\@author
\makeatother

\newcommand{\togglepaper}[1][0]{
   \bibliography{../localbibliography}
   \papernote{\scriptsize\normalfont
     \theauthor.
     \titleTemp.
     To appear in:
     E. Di Tor \& Herr Rausgeberin (ed.).
     Booktitle in localcommands.tex.
     Berlin: Language Science Press. [preliminary page numbering]
   }
   \pagenumbering{roman}
   \setcounter{chapter}{#1}
   \addtocounter{chapter}{-1}
}
%
\newbool{bookcompile}
\booltrue{bookcompile}
\newcommand{\bookorchapter}[2]{\ifbool{bookcompile}{#1}{#2}}
%
% %\newcommand{\orcid}[1]{}
\graphicspath{ {./figures/} }
%
\let\eachwordone=\itshape


% Seongha
% \newcommand{\koreex}[1]{{\fontspec{Baekmuk Batang}{#1}}}

% Kutida
\newfontfamily\thaifont{Noto Sans Thai} 
\newcommand{\textthai}[1]{{\thaifont #1}} % für Thai


\newfontfamily\chinesefont{Noto Sans CJK SC}
\newcommand{\textchinese}[1]{{\chinesefont #1}} % Für Chinesisch

% wiemer
\def\hyph{-\penalty0\hskip0pt\relax}
% \definechangesauthor[color=NavyBlue]{SH}



\newcolumntype{K}{>{\raggedright\arraybackslash}p{0.6\textwidth}}  % Right column with 0.7 width


% \newfontfamily\krn[Scale=MatchLowercase,BoldFont=SourceHanSerif-Bold.otf]{SourceHanSerif-Regular.otf}
% \AdditionalFontImprint{Source Han Serif KO}
% \XeTeXlinebreaklocale 'ko'
% \renewcommand{\krn}{}



% Rhee
\newfontfamily\koreanfont{Noto Sans CJK KR} 
\newcommand{\textkorean}[1]{{\koreanfont #1}} % für Koreanisch

\renewcommand{\REF}[2][]{(\ref{#2}#1)}
	%% hyphenation points for line breaks
%% Normally, automatic hyphenation in LaTeX is very good
%% If a word is mis-hyphenated, add it to this file
%%
%% add information to TeX file before \begin{document} with:
%% %% hyphenation points for line breaks
%% Normally, automatic hyphenation in LaTeX is very good
%% If a word is mis-hyphenated, add it to this file
%%
%% add information to TeX file before \begin{document} with:
%% %% hyphenation points for line breaks
%% Normally, automatic hyphenation in LaTeX is very good
%% If a word is mis-hyphenated, add it to this file
%%
%% add information to TeX file before \begin{document} with:
%% \include{localhyphenation}

\hyphenation{
    con-tri-bute
    par-a-digm
    pe-ri-phra-ses
    achieve-ment
    accomp-lish-ment
    We-li-nchen-kang-ci-kok
    cross-ling-u-is-tic
    rec-ord-s
    trun-ca-ted
    de-ri-va-tion-al
    stem-de-ri-va-tion-al
    pre-sent
    in-flu-en-ced
    sys-tem
    ques-tions
    Bor-kovs-kij
    Da-ny-len-ko
    e-qui-va-lents
    re-cog-nized
    In-ter-fe-renz-schicht
    Aus-ei-nan-der-se-tz-ung
    Früh-neu-hoch-deutsch-kor-pus
    Ak-tions-art
    Smir-no-va
    Whi-le
    Ar-chive
    bi-ginnan
    Rei-chen-bach
    ger-ma-nis-ti-sche
    Re-fe-renz-kor-pus
    le-xi-ka-lisch-en
    Alt-hoch-deutsche
    Li-te-ra-tur
    dia-chron
    Dia-lek-to-lo-gie
}

\hyphenation{
    con-tri-bute
    par-a-digm
    pe-ri-phra-ses
    achieve-ment
    accomp-lish-ment
    We-li-nchen-kang-ci-kok
    cross-ling-u-is-tic
    rec-ord-s
    trun-ca-ted
    de-ri-va-tion-al
    stem-de-ri-va-tion-al
    pre-sent
    in-flu-en-ced
    sys-tem
    ques-tions
    Bor-kovs-kij
    Da-ny-len-ko
    e-qui-va-lents
    re-cog-nized
    In-ter-fe-renz-schicht
    Aus-ei-nan-der-se-tz-ung
    Früh-neu-hoch-deutsch-kor-pus
    Ak-tions-art
    Smir-no-va
    Whi-le
    Ar-chive
    bi-ginnan
    Rei-chen-bach
    ger-ma-nis-ti-sche
    Re-fe-renz-kor-pus
    le-xi-ka-lisch-en
    Alt-hoch-deutsche
    Li-te-ra-tur
    dia-chron
    Dia-lek-to-lo-gie
}

\hyphenation{
    con-tri-bute
    par-a-digm
    pe-ri-phra-ses
    achieve-ment
    accomp-lish-ment
    We-li-nchen-kang-ci-kok
    cross-ling-u-is-tic
    rec-ord-s
    trun-ca-ted
    de-ri-va-tion-al
    stem-de-ri-va-tion-al
    pre-sent
    in-flu-en-ced
    sys-tem
    ques-tions
    Bor-kovs-kij
    Da-ny-len-ko
    e-qui-va-lents
    re-cog-nized
    In-ter-fe-renz-schicht
    Aus-ei-nan-der-se-tz-ung
    Früh-neu-hoch-deutsch-kor-pus
    Ak-tions-art
    Smir-no-va
    Whi-le
    Ar-chive
    bi-ginnan
    Rei-chen-bach
    ger-ma-nis-ti-sche
    Re-fe-renz-kor-pus
    le-xi-ka-lisch-en
    Alt-hoch-deutsche
    Li-te-ra-tur
    dia-chron
    Dia-lek-to-lo-gie
}
	\boolfalse{bookcompile}
	\togglepaper[9]%%chapternumber
}{}

\begin{document}
	\maketitle

	\newpage
	
	\section{Introduction}\label{sec:meyer:intro}
	As an Indo-European language with a respectable number of linguistic idiosyncrasies, both inherited and borrowed,\footnote{\citet[iv]{Olsen1999} famously described Armenian as the ``horror chamber of historical phonology'' owing to its treatment of the cluster PIE *dw- as Arm. \textit{erk-} (e.g. in the word for ‘two’, Arm. \textit{erku} < *dwóh\textsubscript{1}) or the metathesis in the PIE cluster *-br- to Arm. \textit{-łb-} (e.g. in \textit{ełbayr} ‘brother’ as compared to Gk. \textit{pʰrā́tēr}, Lat. \textit{frāter} < PIE *bʰréh\textsubscript{2}tḗr). For the question of borrowing, see the various loans – lexical, phonological, morphological, and syntactic – discussed in \citet{Meyer2023}.} the oldest variety of Armenian, that is Classical or Old Armenian, has received a considerable amount of attention from historical and comparative linguists as concerns most domains, viz. phonetics and phonology, morphology, syntax, semantics, etc.\footnote{See, e.g., the volumes by \citet{Meillet1936}, \citet{Godel1975}, \citet{Klingenschmitt1982} or \citet{Schmitt2007}, which outline the development of Old Armenian from Proto-Indo-European in detail, and the etymological dictionary of \citet{Martirosyan2010}. The most recent, if compendious, discussion can be found in \citet[1028–1165]{KleinJosephFritz2017}'s \textit{Handbook of Comparative and Historical Indo-European Linguistics}.} at least since the latter part of the 19th century; the linguistic study of modern varieties has similarly not been neglected. By far the longest phase of the history of the Armenian language, however, remains regrettably understudied: the Middle Armenian period, stretching from, at the latest, the 12th century to the 18th century and beyond, has received comparatively little attention whilst, at the same time, being the period of greatest diversification.
	
	In turn, this research gap means that perspectives on the diachrony of Armenian as a whole are limited, at least to some extent, by the \textit{tabula rasa} that is the Middle Armenian phase. In order to help remedy this situation, this chapter seeks to delineate and explain the diachronic development of the future from proto-Armenian to its modern varieties. Putting its history in a typological and comparative context, the chapter makes two principal claims: firstly, that the synthetic \textit{k-}future in Modern Eastern Armenian (\textsc{mea}) is the result of two separate instances of hypoanalysis in the Middle Armenian phase; and secondly, that the aspectual diversification of future reference by means of periphrases is the result of analogical paradigmatisation – even though perhaps overstated by current grammars when compared to actual use patterns.
	
	\sectref{sec:meyer:carm} outlines the point of departure for the development of future expression in Armenian – that is, its absence in Old Armenian – and compares it to the situation in other early Indo-European languages. Thereafter, the two main changes in the Middle Armenian period, that is first the creation of a periphrastic present and its later reanalysis as a future, will be discussed in \sectref{sec:meyer:marm}. \sectref{sec:meyer:modarm}, in turn, outlines the current situation in Modern Eastern Armenian and provides some insights into the aspectual diversification in the latter variety.
	
	
	\section{The future in Proto- and Classical Armenian}\label{sec:meyer:carm}
	The oldest attested variety of Armenian, sources for which date back to the 5th century CE,\footnote{While the earliest epigraphic evidence of Armenian dates to the 5th century CE \citep{Stone19901}, the literary attestations post-date this period by far, with the earliest manuscript dating to the 9th century.} does not conceive of the future as a grammatical category. While it distinguishes a present tense from two distinct past tenses – the imperfective preterite and the perfective aorist – and a stative-resultative periphrastic perfect,\footnote{For a discussion of the diachronic semantics of the perfect, cf. \citet{Meyer2024a}.} future reference can only be expressed by means of polysemous forms. As examples (\ref{ex:meyer:carm1}–\ref{ex:meyer:carm2}) and (\ref{ex:meyer:carm3}–\ref{ex:meyer:carm4}) illustrate, this is done either by means of the subjunctive, which has a variety of other, related functions,\footnote{They can, for instance, express doubt or uncertainty in a variety of subordinate clauses \citep[200, 205--207]{Jensen1959}.} or by means of a periphrasis composed of the deontic verbal adjective in \textit{-loc‘} and a form of the copula.\footnote{The etymology of this verbal adjective is not clearly understood. \citet[129]{Godel1975} plausibly proposes that it is composed of the formans *-lo-, which is also used to form the infinitive and past participle, and a second, derivative suffix *-sko- as found in Germanic and Slavic; the functional parallels are, however, less clear.}
	
	\ea \label{ex:meyer:carm1}
	\gll erkink‘ ew erkir anc‘c‘en\\
	heaven.\textsc{nom.pl} and earth.\textsc{nom.sg} pass.\textsc{3pl.sbjv}\\
	\glt ‘Heaven and earth shall pass’ (Mt. 24:35)
	\z \il{Old Armenian}
	
	\ea \label{ex:meyer:carm2}
	\gll ew es ǝnd jez ibrew z=mi i jēnǰ xndrec‘ic‘ z=jer ōgut=n\\
	\textsc{conj} \textsc{1sg.nom} with \textsc{2pl.acc} like \textsc{obj=one.acc.sg} among \textsc{2pl.abl} seek.\textsc{1sg.aor.sbjv} \textsc{obj=2pl.poss} advantage.\textsc{acc.sg=det}\\
	\glt ‘And I, too, like one of you, shall seek your advantage’ (Agat‘angełos CI.9; \citealt{Matenagirk2})
	\z \il{Old Armenian}
	
	\ea \label{ex:meyer:carm3}
	\gll mi omn i jēnǰ matneloc‘ ē z=is\\
	one \textsc{indf.nom.sg} from \textsc{2pl.abl} betray.\textsc{vbadj} be.\textsc{3sg.prs} \textsc{obj=1sg.acc}\\
	\glt ‘One among you will betray me’ (Mk. 14:18)
	\z \il{Old Armenian}
	
	\ea \label{ex:meyer:carm4}
	\gll zi ew duk‘ ibrew ołormut‘eamb K‘ristosi arjakeloc‘ ek‘ i kapanoc‘=d, ew mek‘ … gnaloc‘ emk‘ i bnik ašxarh=n mer\\
	for also \textsc{2pl.nom} like mercy.\textsc{ins.sg} Christ.\textsc{gen.sg} release.\textsc{vbadj} be.\textsc{2pl.prs} from chain.\textsc{abl.pl=det} and \textsc{1pl.nom} {} go.\textsc{vbadj} be.\textsc{1pl.prs} to natural land.\textsc{acc.sg=det} \textsc{1pl.poss}\\
	\glt ‘For you, by the mercy of Christ, shall be released from your chains, and we … shall return to our homeland’ (Łazar P‘arpec‘i III.50.1; \citealt{Matenagirk2})
	\z
	\il{Old Armenian} %add "Latin" to language index for this page
	
	While this polysemy between deontic or potential uses of the forms on the one hand and future uses on the other has the disadvantage of allowing for a certain degree of ambiguity, it is neither typologically, historically, nor comparatively uncommon.
	
	No future category can be reconstructed for Proto-Indo-European, for instance. Those early languages like Sanskrit, Greek or Latin which do attest a synthetic future and thus a distinct grammatical category have developed it individually by reanalysis of a formerly subjunctival or desiderative form or by means of univerbation of a previous periphrastic construction.\footnote{For a discussion of the desiderative or subjunctival origin of the Greek sigmatic future, 
    see \citet[441--452]{Willi2018}. The Latin \textit{b-}future appears to be a univerbated form of the verbal root and the aorist subjunctive of *bʰuH- ‘to become’ \citep[415]{Weiss2009}; the original form and function of the verbal root is, however, as yet unclear.} Others, like Armenian, rely on the subjunctive and its polysemy to express futurity; one such example is Tokharian (\ref{ex:meyer:tokh}).\footnote{Cf. \citet[265–266]{Malzahn2010}.}
	
	\ea \label{ex:meyer:tokh}
	\gll näṣ nu ce(smä)k āyäntu p\textsubscript{u}kāk puskāsyo kaśal malkam=äm\\
	\textsc{1sg.nom.m} and \textsc{dem.acc.pl.m.emph} bone.\textsc{acc.pl.m} all.\textsc{perl} sinew.\textsc{ins.pl.m} together join.\textsc{1sg.sbjv.act=obj.pl}\\
	\glt ‘And I will join the bones completely together with the sinews’ (Tokharian A; A 11 b5–6)
	\z
	\il{Tokharian}

	\largerpage[-1.5]
	
	Next to the subjunctive, Tokharian, among other Indo-European languages, can employ the present tense  effectively as a non-past, thus encompassing future events; at times, adverbs and serial verb constructions are used to elucidate this futurity and distinguish it from the present,\footnote{Hittite, for instance, can use \textit{kāša} ‘at the point of’ for the immediate future or \textsc{\textit{urram šeram}} ‘in the future’ as in (\ref{ex:meyer:hitt}) below \citep[324]{HoffnerMelchert2008}. Gothic can use \textit{wairþan} ‘to get to be’ as a prophetic future, or periphrases with \textit{du-ginnan} ‘to begin’, \textit{haban} ‘to have’, or \textit{skulan} ‘to owe, be obliged to’ \citep{Kleyner2015}. Old Church Slavonic equally can use \textit{xotĕti} ‘to want’ or \textit{imamǐ} ‘to have’ to express futurity \citep{Lunt2001}.} without being sufficiently widespread or grammaticalised, however, as to count as a separate tense category.
	
	\ea \label{ex:meyer:hitt}
	\gll {\textsc{urram} \textsc{šeram}} kuiš ammuk \textsc{egir}-anda \textsc{lugal}-uš kišari\\
	in-the-future \textsc{rel.nom} \textsc{1sg.dat} after king-\textsc{nom} become.\textsc{3sg.prs.mid}\\
	\glt ‘Whoever shall become king after me’ (Hittite; KBo 3.1 ii 40)
	\z
	\il{Hittite}
	
	\ea \label{ex:meyer:goth}
	\gll  saei bigitiþ saiwala seina, fraqisteiþ izai\\
	\textsc{rel.nom.sg.m} find.\textsc{3sg.prs.ind} soul.\textsc{acc.sg.f} \textsc{3sg.poss.acc.f} lose.\textsc{3sg.prs.ind} \textsc{3sg.dat.f}\\
	\glt ‘Whoever finds his soul, will lose it’ (Gothic; Mt. 10:39 [Argenteus])
	\z
	\il{Gothic}

	\ea \label{ex:meyer:ocs}
	\gll nebo i zemlę mimo idetǔ\\
	heaven.\textsc{nom.sg} and earth.\textsc{nom.pl} past go.\textsc{3sg.prs}\\
	\glt ‘Heaven and earth shall pass’ (Old Church Slavonic; Mt. 24:35 [Marianus])
	\z
	\il{Old Church Slavonic}
	
	The above examples illustrate such uses for Hittite (\ref{ex:meyer:hitt}) as a language that separated early from the Indo-European core;\footnote{On the position of the Anatolian languages, see most recently \citep{Kloekhorst2022}.} and for Gothic (\ref{ex:meyer:goth}) and Old Church Slavonic (\ref{ex:meyer:ocs}) as languages with a later attestation.\footnote{On Gothic, cf. \citet[176 n. 1]{Miller2019} with literature; for Old Church Slavonic, cf. \citet[154–155]{Lunt2001}.}
	
	Even beyond the context of Indo-European languages alone, synthetic futures are by no means commonplace. In \textcite{DahlVelupillai2013}'s typological study, just under 50\% of the 222 studied languages exhibited an inflectional future category; for \citet[156]{Bybee1985}, it was 44\% in a sample of 50. Historically speaking, however, even these inflectional futures are often the result of grammaticalised, and thus univerbated, periphrases originally expressing volition, obligation, or motion \citep[159]{BybeePerkinsPagliuca1994};\footnote{For an overview of such developments in other language families, see, e.g., \citet{AbdelHafiz2017} for Arabic and the Nile-Nubian languages; for Yucatecan Maya, cf. \citet[213–229]{Lehmann2017}; 
	for Bantu, cf. \citet{Batibo2005} 
	and \citet{Botne1998} for a future in Central Bantu based on a verb meaning ‘to say’.} the Romance languages demonstrate neatly that this process of grammaticalisation can by cyclical: the development of the deontic collocation of the type \textit{cantāre habeō} ‘I have to sing’ into a future expression in Late Latin overlapped with the loss of the synthetic future (type \textit{cantābō}), resulting in the eventual grammaticalisation and standardisation of the (previously periphrastic) forms as the new ‘standard’ future of the type Fr. \textit{chanterai}, It. \textit{canterò}, Sp. \textit{cantaré} ‘I shall sing’; these, in turn, compete in modern varieties with motion-related periphrases of the type Fr. \textit{aller chanter}, Sp. \textit{ir a cantar} ‘going to sing’.\footnote{Italian does not have such a periphrasis. Similarly, both German and English have a fully grammaticalised periphrastic future, but only English has developed a secondary, motion-based future in \textit{going to}.}\il{Latin, Spanish, Italian, French}
	
	Old Armenian therefore is, so to speak, in good company in not (yet) having a morphologised future, but equally exhibits periphrases which might potentially undergo such a process. As it turns out, however, this is not the pathway that the Armenian language will follow; as the next section explains, neither the (synchronic) subjunctive nor the deontic periphrasis will serve as the basis for a later future tense. Rather, in Middle Armenian, two cycles of grammaticalisation are required before the crystallisation of a separate future category. These are not dissimilar to the Romance cycles just mentioned, but centre around the present rather than the future tense, with the latter's creation being a by-product.
	
	
	\section{Developments in Middle Armenian}\label{sec:meyer:marm}
	
	\largerpage[1.5]

	In the more or less six centuries of the Middle Armenian period, the Armenian language has undergone significant changes, both in terms of its phonology and morphology as well as its general unity. Owing to the dearth of research in this area, some of these processes can only be judged by their outcomes at present, rather than by a detailed analysis of the changes itself. In terms of diversity, at the beginning of the 20th century, \citet{Acarean1911} describes 41 distinct varieties in three groups.\footnote{An English translation with a commentary of this magnum opus is due to be published \citep{DolatianAdjarianfthc}.} Not all of these varieties have undergone the same changes: the western so-called \textit{gə}-varieties, for instance, have lost the threefold distinction of plosives in voiced (/b/, /d/, /g/), voiceless (/p/, /t/, /k/), and unvoiced aspirated (/pʰ/, /tʰ/, /kʰ/) attested for Old Armenian and, with some phonetic variation, retained in the Eastern, \textit{um-}varieties (including Modern Eastern Armenian, the standard literary variety); the Western varieties only retain a distinction between voiced and voiceless aspirate.\footnote{The old voiceless aspirate plosives retain their articulation, whereas the former voiced plosives fusion with them (/b/, /d/, /g/ > /pʰ/, /tʰ/, /kʰ/); the old unvoiced plosives are voiced (/p/, /t/, /k/ > /b/, /d/, /g/).}
	
	This phonological diversification together with other changes seems to belie the notion of Middle Armenian as a unified whole representing all of high- and late-medieval and early modern Armenian. This complex situation is not helped by the fact that research on Middle Armenian to date has focussed largely on phonological, lexical, and dialectal questions as well as language contact,\footnote{For an overview with plentiful literature, cf. \citet{vanLintfthc}; the question of Middle Armenian and its dialectal variety is treated, amongst others, in \citet{Jahukyan1993},  \citet{Mkrtcyan2018},  \citet{Muradyan1975},  \citet{Weitenberg1996} and \citet{Weitenberg2017}.} with very little in terms of morphological or morphosyntactic research.\footnote{Notable exceptions are \citet{Antosyan1972}, \citet{Antosyan1975}; for more general linguistic grammatical overviews, cf. \citet{Alayan1972}, \citet{Hovsepyan1975}, \citet{Jahukyan1964}, \citet{Jahukyan1975} and  \citet{Lazaryan1960}.} As a result, the standard work of reference remains \citet{Karst1901}, which, alas, gives only limited detail on verbal morphosyntax in general, is not primarily concerned with diachronic developments, and in general only discusses one particular variety of Middle Armenian -- the variety of Cilicia. For the present purpose, Middle Armenian is thus treated as a whole, all the while acknowledging that this is a simplification.
	
	To properly understand the diachrony of future expression in Middle Armenian, therefore, the most promising approach would seem to be triangulation: to take what is known for certain in Old Armenian and the modern varieties in order to reconstruct the most plausible trajectory of the changes involved, including eventual forks in the path. For that purpose, it is worth observing changes not only in expressions of futurity but also, as already alluded to above, in the present tense system.
	
	\tabref{tab:meyer:prsfut1} summarises the situation: Old Armenian uses a simple synthetic indicative form to express the present tense, whereas the future is marked by (aorist) subjunctive forms or a deontic periphrasis. The modern standard varieties (Modern Eastern Armenian, \textsc{\textsc{mea}}, and Modern Western Armenian, \textsc{\textsc{mwa}}) use different formations to express future and present, respectively.\footnote{The more fine-grained, mostly phonological differences that are found in the present and future formations in the individual dialects will not be discussed here.} In \textsc{\textsc{mea}}, the present is formed periphrastically with a verbal adjective in \textit{-um} and a finite form of the copula; the (unmarked) future is formed by adding the future marking morph \textit{k-} to what is historically the Old Armenian present tense and synchronically the subjunctive. In \textsc{\textsc{mwa}}, by contrast, it is the present which is formed by adding the marker \textit{kə} (or \textit{k'}) to the historical present / synchronic subjunctive;\footnote{These morphs are pronounced as /g(ə)/ in \textsc{\textsc{mwa}}; in this paper, however, Armenian is consistently transcribed according to Hübschmann-Meillet-Benveniste as the convention used in most subject-specific publications and journals, e.g. the \textit{Revue des Études Arméniennes}.} the future is formed by means of a historically deontic periphrasis with \textit{piti} (/bidi/) and the historical present / synchronic subjunctive.
	
	\begin{table}
		\caption{Future and present in Old and Modern Armenian}
		\label{tab:meyer:prsfut1}
		\begin{tabularx}{1\textwidth}{l llcll r}
			\lsptoprule    & Old                                       &            &  Middle  &            & Modern              &     \\
			               & Armenian                                  &            & Armenian &            & Armenian            &     \\ \midrule
			               &                                           &            &          &            & \textit{berum em}   & \textsc{mea} \\
			\textsc{prs}   & \textit{berem}                            &            &          &            & ‘I (am) carry(ing)’ &     \\
			               & ‘I carry’                                 & $\searrow$ &          & $\nearrow$ & \textit{kə berem}   & \textsc{mwa} \\
			               &                                           &            &          &            & ‘I (am) carry(ing)’ &     \\
			               &                                           &            &    ?     &            &                     &     \\
			               &                                           &            &          &            & \textit{kberem}     & \textsc{mea} \\
			\textsc{fut}   & {[\textit{beric‘} / \textit{bereloc‘ em}]} & $\nearrow$ &          & $\searrow$ & ‘I shall carry’     &     \\
			               & {[‘I shall carry’]}                        &            &          &            & \textit{piti berem} & \textsc{mwa} \\
			               &                                           &            &          &            & ‘I shall carry’     &     \\
			\lspbottomrule &                                           &            &          &            &                     &
		\end{tabularx}
	\end{table}
	
	These various formations allow for a certain number of observations of commonalities and differences to be made; they are, in no particular order:
	\begin{itemize}
		\item the Old Armenian future periphrases do not have modern cognates;
		\item the Old Armenian present formation is continued in the modern varieties (as a subjunctive) and involved in both present (\textsc{mwa}) and future formation (\textsc{mea}, \textsc{mwa});
		\item the marker \textit{k- / kə} is used in both modern varieties, but for different purposes – the future in \textsc{mea}, the present in \textsc{mwa};
		\item both modern varieties have grammaticalised one tense on the basis of a periphrasis not found in that form in the other variety (\textsc{mea} present; \textsc{mwa} future).
	\end{itemize}
	
	In terms of diachronic expectations, it therefore stands to reason that, in the first place, the Old Armenian futuric periphrases fell out of use (or, perhaps, were never sufficiently frequent to cross the grammaticalisation threshold) and the subjunctive was lost. Secondly, the use of forms resembling the Old Armenian present indicative as subjunctives in the modern varieties suggests a shift in that direction in Middle Armenian; such developments are, of course, well-attested and typologically unproblematic \citep[cf.][]{Haspelmath1998a}. Concomitant with the reanalysis of the old present indicative as an innovative subjunctive, a new present indicative must have arisen to motivate this change to begin with. This innovative present must have involved the predecessor of the marker \textit{k- / kə}, on the assumption of economy of change.\footnote{If there was a \textit{k-}present in Middle Armenian, only one change is required to explain the current distribution: the Middle Armenian present becomes an \textsc{mea} future. Otherwise, if a different grammatical value is attributed to this form, multiple changes, namely for both modern varieties, would be required.} Western Armenian varieties would thus be expected to have retained this constellation, and innovated a deontic future periphrasis, whereas further change ensued in the Eastern Armenian varieties. There, a second round of reanalysis must have resulted in the reinterpretation of the \textit{k-}present as a future under competition from a new present formation.
	
	These are, of course, only working hypotheses based on a reasoned comparison of the earliest and recent varieties of Armenian. They could and should, in the long run, be verified and (dis-)proved on the basis of detailed corpus studies. In the absence of such a corpus, however, the following sections provide relevant illustrative examples from Middle Armenian, which have to suffice in order to show that the progress outlined above is, at the very least, plausible and to provide more details of the mechanisms of change.
	
	\subsection{Early Middle Armenian}\label{sec:meyer:emarm}
	
	The first question arising concerns the fate of the Old Armenian subjunctive. As a result of its relatively small functional load besides the expression of the future and due to a number of morphophonological idiosyncrasies,\footnote{Cf. \citet[299]{Karst1901}; \citet{Weitenberg1993}; \citet{Vaux1995}.}
	both the present and aorist forms fell out of use already towards the end of the Old Armenian period.\footnote{A reviewer inquires whether the loss of the old subjunctive might be the result of the advent of the \textit{ku}-present and the subsequent reanalysis of the old present as the subjunctive. Without a detailed corpus study, this cannot be excluded \textit{a priori}. Two issues may speak against it, however: the lack of co-occurrence of all three forms (old subjunctive, old present, and \textit{ku-}presents) in early Middle Armenian texts as might be expected in a transitional period; and the complete loss of the old subjunctive instead of its relegation to a more specific function, e.g. as a future.} Their functions are assumed by the present indicative, as witnessed by the virtual identity (barring phonological change) of modern Armenian present subjunctives and old present indicatives;\footnote{Cf. \citet{Acarean1911} \textit{passim} and the relevant commentaries by \citet{DolatianAdjarianfthc}.} this inevitably increased the functional load of that form significantly. A clear indication of this replacement can be found in the expression of commands or exhortations, which in Old Armenian require the subjunctive (\ref{ex:meyer:osubj1}–\ref{ex:meyer:osubj2}) but in Middle Armenian are taken over by the present indicative and a particle (\ref{ex:meyer:subj1}–\ref{ex:meyer:subj2}).
	
	\ea \label{ex:meyer:osubj1}
	\gll ari gnasc‘uk‘ aṙ kaysr=n\\
	valiant go.\textsc{2pl.aor.sbjv} to emperor.\textsc{acc.sg=det}\\
	\glt ‘Let us go valiantly to the Emperor’ (P‘awstos Buzand IV.10; \citealt{Matenagirk1})
	\z \il{Old Armenian}
	
	\ea \label{ex:meyer:osubj2}
	\gll iwrak‘anč‘iwrok‘ z=iwr kin kalc‘i\\
	everyone \textsc{obj=3sg.poss} woman.\textsc{acc.sg} have.\textsc{3sg.aor.sbjv}\\
	\glt ‘Let everyone have his own wife’ (1Cor. 7:2)
	\z \il{Old Armenian}
	
	\ea \label{ex:meyer:subj1}
	\gll t‘ē jmeṙn lini, tak‘ ǰur t‘oł xmē\\
	if winter.\textsc{nom.sg} become.\textsc{3sg.prs} warm water.\textsc{acc.sg} \textsc{ptc} drink.\textsc{3sg.prs}\\
	\glt ‘If it is winter, let him drink warm water.’ (Mxit‘ar Herac‘i, \textit{J̌ermanc‘ mxitarut‘iwn} 20.14)
	\z \il{Middle Armenian}
	
	\ea \label{ex:meyer:subj2}
	\gll t‘oł mašim i y=ays hog=s ew xist hənanam\\
	\textsc{ptc} be-consumed.\textsc{1sg.prs} into into=\textsc{det} care=\textsc{det} and firm grow-old.\textsc{1sg.prs}\\
	\glt ‘Let me be consumed by this care and grow old [in] firm [belief]’ (Frik, \textit{Diwan} 38.86)
	\z \il{Middle Armenian}
	
	The exhortation or command signified by the subjunctive forms in Old Armenian can no longer be expressed by such forms in Middle Armenian, as they are identical to the indicative;\footnote{\label{fn:meyer:1}\citet[299–300]{Karst1901} mentions the existence of a subjunctive periphrasis consisting of the indicative forms of the verb \textit{kam} ‘remain’, the conjunction \textit{u} ‘and’, and a finite lexical verb (e.g. \textit{et‘ē kenay u gay} ‘if he comes’), but this periphrasis is short-lived and never fully grammaticalised.} the necessary functional clarification is added by the particle \textit{t‘oł}.\footnote{This particle derives from \textit{t‘ołum} ‘to let, allow’, of which it is a, by then fossilised, imperative. Such a construction is already attested in Old Armenian, e.g. \textit{t‘oł ekesc‘ē dma uraxut‘iwn} ‘let happiness come for him’ (Agat‘. 105); note, however, that here a subjunctive is still required.} In other instances yet, old structures or collocations are lost owing to the increased formal ambiguity; such is the case, for instance, in purposes clauses introduced by \textit{zi} ‘so that’, where Old Armenian requires a subjunctive (\ref{ex:meyer:osubj3}). In Middle Armenian, this use of \textit{zi} is no longer possible.\footnote{The conjunction remains in active use, however, in its causal function, meaning ‘for’.}
	
	\ea \label{ex:meyer:osubj3}
	\gll {na ew} episkoposac‘=n žołoveal, pałatans matuc‘anēin zi mi anp‘oyt‘ arasc‘ē z=korstenē iwroy vičaki=n\\	
	even bishop.\textsc{gen.pl=det} gather.\textsc{cvb}  suplication.\textsc{acc.pl} offer.\textsc{3pl.pst} so-that \textsc{neg} careless do.\textsc{3sg.aor.sbjv} about=loss.\textsc{abl.sg} \textsc{3sg.poss.gen} domain.\textsc{gen.sg=det}\\
	\glt ‘Even the bishops gathered and sent supplications lest he act indifferent concerning the loss of his domain’ (Movsēs Xorenac‘i III.29; \citealt{Matenagirk2})
	\z \il{Old Armenian}

	Next to these subjunctival functions \textit{sensu stricto}, the old present also occasionally takes on futuric functions, as indicated by (\ref{ex:meyer:mfut1}).
	
	\ea \label{ex:meyer:mfut1}
	\gll p‘aṙk‘ Astucoy ver aṙak‘em ew biwr hazar het gohanam, or inj eret mitk‘ u hogi, or z=ays baners ku səgam.\\
	glory god.\textsc{gen.sg} up send.\textsc{1sg.ind} and hundred thousand time thank.\textsc{1sg.ind} \textsc{comp} \textsc{1sg.dat} give.\textsc{3sg.aor} mind and soul \textsc{rel.nom.sg} \textsc{obj=dem.acc.sg} thing.\textsc{acc.pl} \textsc{prog.ind} mourn.\textsc{1sg}\\
	\glt ‘I shall send up high the glory of God and thank [him] a hundred thousand times | that he gave me a mind and a soul, who I am lamenting these things’ (Frik, \textit{Diwan} 2.46–7)
	\z \il{Middle Armenian}
	
	This functional overload – one form for present indicative, subjunctive, and future – proved to be an ideal breeding ground for other, innovative forms to intervene and take on specific functions from the old present. This is not to suggest, of course, that there is necessarily a causal relationship between functional overload and the \textit{creation} of new forms in general. In the case of Middle Armenian, the innovation comes in the guise of the particle \textit{ku}, likely originally a progressive marker.\footnote{\label{fn:meyer:2} \citet{Vaux1995} argues convincingly for a progressive reading of this particle post-standardisation, since it does not occur with a number of high-frequency stative verbs; \citet[307–309]{Karst1901} rightly remarks, however, that the (tentative) derivation of \textit{ku} from \textit{kay u} ‘remains and’ and its general distribution strongly suggest that the particle originally signified mainly the indicative.} This periphrasis does not take on subjunctival roles,\footnote{But see fn. \ref{fn:meyer:1} above and compare the origin of the present particle as mentioned in \ref{fn:meyer:2}, a certain likeness between which cannot be denied.} but expresses mainly the present (\ref{ex:meyer:mprs1}–\ref{ex:meyer:mprs2}); from Mxit‘ar Herac‘i (12th c.) onward it is indeed the standard expression of the present tense.
	
	\ea \label{ex:meyer:mprs1}
	\gll menk‘ ekel enk‘ u ku hrawirenk‘ z=k‘ez or gas sur i darpas=n\\
	\textsc{1pl.nom} come.\textsc{ptcp} be.\textsc{1pl.prs} and \textsc{prog.ind} invite.\textsc{1pl.prs} \textsc{obj=2sg.acc} \textsc{comp} go.\textsc{2sg.prs.sbjv} immediately to court.\textsc{acc.sg=def}\\
	\glt ‘We have come and invite you to go to court immediately’ (Smbat Sparapet, \textit{Antiok‘i asizner} 13.13ff.)
	\z \il{Middle Armenian}
	
	\ea \label{ex:meyer:mprs2}
	\gll ew erb imanas or ōr k‘an z=ōr erak=n ku tkaranay, apa z=hac‘=n i gini=n t‘ac‘, ew tur\\
	and when(ever) notice.\textsc{2sg.sbjv} \textsc{comp} day after \textsc{obj}=day vein.\textsc{nom.sg=det} \textsc{prog.ind} weaken.\textsc{3sg.prs} then \textsc{obj=}bread.\textsc{acc.sg=det} in wine.\textsc{acc.sg=det} soak.\textsc{2sg.aor.imv} and give.\textsc{2sg.aor.imv}\\
	\glt ‘And when(ever) you notice that the vein weakens day after day, soak the bread in wine and give [it to him]’ (Mxit‘ar Herac‘i, \textit{J̌ermanc‘ mxitarut‘iwn} 20.17)
	\z \il{Middle Armenian}

	\largerpage[1]
	
	Even if its origins (and later use in \textsc{mwa}) suggest a progressive reading, the periphrasis with \textit{ku} at least in early Middle Armenian can also express futurity and should thus, perhaps, be interpreted as a more general non-past (\ref{ex:meyer:mfut3}); another, more specific future periphrases of a voluntative nature (\textit{kamim} + \textsc{inf}) exists as well, but has not been standardised at this time (\ref{ex:meyer:mfut4}).\footnote{The deontic periphrasis behind the later \textsc{mwa} future (\textit{piti} + \textsc{sbjv}) exists already in Middle Armenian, but there still has only a deontic meaning \citep[306]{Karst1901}.}
	
	\ea \label{ex:meyer:mfut3}
	\gll es ku asem u ku c‘uc‘ənem alani u yerewan\\
	\textsc{1sg.nom} \textsc{prog.prs.ind} say.\textsc{1sg} and \textsc{prog.prs.ind} demonstrate.\textsc{1sg} clearly and evidently\\
	\glt ‘I will speak [at court] and will demonstrate clearly and evidently…’ (Smbat Sparapet, \textit{Antiok‘i asizner} 27.11)
	\z \il{Middle Armenian}
	
	\ea \label{ex:meyer:mfut4}
	\gll ew mek=n i šoyt ew i yuš žołovel=n lini niwt‘i=n, or borbosel kami\\
	and one.\textsc{nom.sg=det} in quick and in slow gather.\textsc{inf=det} become.\textsc{3sg.prs} substance.\textsc{gen.sg=det} \textsc{rel.nom.sg} mould.\textsc{inf} wish.\textsc{3sg.prs}\\
	\glt ‘One [reason] is the gathering, quickly or slowly, of a substance, which will grow mouldy’ (Mxit‘ar Herac‘i, \textit{J̌ermanc‘ mxitarut‘iwn} 24.70)
	\z \il{Middle Armenian}
	
	The early Middle Armenian period is, as has been illustrated, a period of transition in morphosyntax, with a certain degree of functional overlap at various points in time between the old synthetic present forms, which gradually lose their role of expressing the present indicative in favour of reflecting the subjunctive and the future, and the new present periphrasis in \textit{ku}, which also can have futuric readings. The specific value of this particle \textit{ku} is difficult to pin down; while \citet[137]{Vaux1995} believably argues for an original progressive meaning, at the latest by the time of the standardisation of this periphrasis, it had lost the specific meaning and only signified indicative.\footnote{This effectively overrides the by then inherently subjunctival function of the finite verb form. A temporal assignation of present, future, or non-past is not straightforwardly possible, since the past indicative is also formed with this particle (e.g. \textit{ku xmēr} ‘he drank’; Smbat Sparapet, \textit{Taregirk‘} 650) and since both present and future events can be expressed as shown above.}
	
	As will become evident from the changes ensuing in later Middle Armenian, it is after the set of developments discussed here that Eastern and Western Armenian varieties stopped developing their verbal system in parallel (or indeed together as one); Modern Western Armenian retains a verbal system closer to the one described thus far, whereas the Eastern varieties undergo further developments, which are outlined in the following section.
	
	\subsection{Later Middle Armenian}
	The developments in this later Middle Armenian phase are most clearly visible from their outcomes in the various varieties of Armenian described by \citet{Acarean1911} and already referenced above. In the Western varieties, a variant of the \textit{ku-}present persists, albeit often coupled with phonetic reduction of the particle \textit{ku} to \textit{k‘} or \textit{kə}.
	
	In the Eastern varieties, which \citet{Acarean1911} refers to as \textit{-um}-varieties because of their innovation of a periphrastic present ending in \textit{-um}, the situation is evidently different. Here, the old \textit{ku-}present has adopted a futuric function, with all present functions being taken over by an innovative, originally locatival periphrasis. Once more, owing to the present unavailability of a corpus of Middle Armenian texts, it is difficult to establish when this shift took place. Its firm establishment in Tiflis-dialect poetry by the middle of the 18th century, however, suggests this as a reasonable \textit{terminus ante quem}, given that poetry – with the metrical constraints imposed upon it – tends to allow for archaising language. When the Georgian-Armenian \textit{ašuł} Sayat‘-Nova uses innovative forms, therefore, it can be assumed that they reflect a firmly established linguistic norm (\ref{ex:meyer:mmea1}–\ref{ex:meyer:mmea2}).
	
	\ea \label{ex:meyer:mmea1}
	\gll T‘e xam hak‘nis, t‘e zar hak‘nis, ku šinis łumaš, nazani.\\
	whether simple wear.\textsc{2sg.sbjv} whether gold wear.\textsc{2sg.sbjv} \textsc{fut} make.\textsc{2sg} fine darling\\
	\glt ‘Whether you wear simple [things] or gold, you make it fine, [my] darling’ (Sayat‘-Nova 26.3; \citealt{Meyer2023d})
	\z \il{Tiflis Armenian}
	
	\ea \label{ex:meyer:mmea2}
	\gll Łalam=ən jeṙin č‘=ē kangnum, mat‘ šinec‘ir nałašk‘ari=n\\
	pen.\textsc{nom.sg=det} hand.\textsc{loc.sg} \textsc{neg=}be.\textsc{3sg.prs} rest.\textsc{ptcp} mate make.\textsc{2sg.aor} artist.\textsc{dat.sg=det}\\	
	\glt ‘The pen does not rest in [his] hand, you have set the artist mate’ (Sayat‘-Nova 26.15)
	\z \il{Tiflis Armenian}
	
	By contrast, at roughly the same time, the poet Pałtasar Dpir, who writes in the Istanbul variety, employs the same \textit{ku}-forms to denote the present tense (\ref{ex:meyer:mmwa1}), while at the same time retaining older, unmarked presents for similar purposes (\ref{ex:meyer:mmwa2}); it is unclear from the context whether the different forms were in free variation, represented an aspectual contrast (e.g. dynamic vs habitual), or whether this particular variety of Middle Armenian is informed by the poet's educational aspirations.\footnote{Pałtasar sought to revitalise Old Armenian as a literary language; it cannot, therefore, be excluded that he should have been influenced by that variety's morphosyntax.}
	
	\ea \label{ex:meyer:mmwa1}
	\gll ōtari=n lezu=n aweli c‘aw ełer, z=is ku tałē, ku tałē\\
	foreign.\textsc{gen.sg=det} tongue.\textsc{nom.sg=det} too pain.\textsc{nom.sg} become.\textsc{3sg.pst} \textsc{obj=1sg.acc} \textsc{ind} burn.\textsc{3sg.prs} \textsc{prs} burn.\textsc{3sg}\\
	
	\glt ‘The foreign tongue became too painful, it burns me, it burns’ (Pałtasar Dpir, \textit{Poems} 1.2)
	\z \il{Istanbul Armenian}
	
	\ea \label{ex:meyer:mmwa2}
	\gll Aregakn cagē y=erknic‘ i y=erkir, \emph{[…]} erb yišem z=k‘ez, linim orpēs z=moxir\\
	sun.\textsc{nom.sg} shine.\textsc{3sg.prs} from=sky.\textsc{abl.pl} on to=earth.\textsc{acc.sg} {} when remember.\textsc{1sg.prs} \textsc{obj=2sg.acc} become.\textsc{1sg.prs} like \textsc{obj=}ash.\textsc{acc.sg}\\
	\glt ‘The sun shines from the sky unto earth, […] when I think of you, I burn up like ashes’ (Pałtasar Dpir, \textit{Poems} 2.16–18)
	\z \il{Istanbul Armenian}
	
	While there is a certain inherent logic in the notion that varieties which undergo further systemic changes should show greater stabilisation and differentiation of the forms that are the basis for said changes — in this instance, the Eastern \textit{-um}-varieties — it must be kept in mind that at least more highly educated Armenian speakers were likely diglossic. Next to their local, spoken varieties existed a more or less standardised (and more or less archaic or archaising) written variety, following Old Armenian morphosyntax and retaining its orthographic conventions, which was still used in the Church and for official documents. In some cases, it therefore remains uncertain to what extent the varietal texts in existence reflect purely the spoken language, and to what extent the diglossic situation has permitted the old, written standard to influence (or be influenced by) the more modern variety.
	
	Returning to the question of Middle Armenian developments in the Western varieties, the question of the future remains to be discussed. As already mentioned, these varieties retain the \textit{ku-} forms in their early Middle Armenian function as presents alone and had to innovate a separate form for the future. The periphrasis that eventually prevailed is built on a historically deontic expression with \textit{piti} (\ref{ex:meyer:mmwa3}).\footnote{The form \textsc{mwa} \textit{piti} is derived from OArm. \textit{pētk‘} ‘necessary’. In Modern Eastern Armenian both \textit{piti} and \textit{petk‘} are still used to express necessity (debitive), e.g  \textsc{mea} \textit{petk‘ ē štapenk‘} ‘we must hurry’; \citet[263–271]{DumTragut2009} notes that, in spoken Eastern Armenian, the form \textit{piti} is more common.}
	
	\ea \label{ex:meyer:mmwa3}
	\gll ayn ordin or kə całrē ir hayr=ə, noyn vičaki=n piti matnui ōr mə.\\
	\textsc{dem} son.\textsc{nom.sg=det} \textsc{rel.nom.sg} \textsc{ind} mock.\textsc{3sg.prs} \textsc{3sg.poss.refl} father.\textsc{acc.sg=det} same.\textsc{nom.sg} fate.\textsc{dat.sg=det} \textsc{fut} face.3sg day \textsc{indf}\\	
	\glt ‘The son who mocks his own father will one day face the same fate.’ (Dawit‘ Xač‘konc‘, \textit{Dun piti əllas baroyagēt, tntesagēt, ōrēnsgēt} (1907), 79)
	\z \il{Modern Western Armenian}
	
	While this example post-dates the Middle Armenian period by some time, it illustrates most clearly the key difference between Modern Western and Eastern Armenian. While the former retains the present innovated in early Middle Armenian and innovates a periphrastic future, the latter undergoes a further development (in later Middle Armenian), retaining \textit{ku-}forms created earlier in their secondary, future function and innovating a present form instead; at a later stage, this future periphrasis is univerbated. In both varieties, the subjunctive is the most stable form, retained in more or less the same form since early Middle Armenian (and itself based on the Old Armenian present indicative). Notably, the subjunctive form is also the basis for some of the innovative periphrases (\textit{ku-}presents and \textit{ku-}futures as well as the \textit{piti}-future). \tabref{tab:meyer:prsfut2} provides a diagrammatic overview of these developments.

	\begin{table}
		\caption{Future and present from Old to Middle and Modern Armenian. Cells shaded in dark grey indicate forms no longer employed in this function after said time period; cells in light grey indicate innovations.}
		\label{tab:meyer:prsfut2}

	\begin{tabularx}{.8\textwidth}{Xllll}
		\lsptoprule       &               & \textsc{prs}                            & \textsc{sbjv}             & \textsc{fut}                            \\ \midrule
		OArm.             &               & \tikzmark{1}berem                       & \cellcolor{gray!60}beric‘ & \cellcolor{gray!60}bereloc‘ em          \\
		                  &               & \tikzmark{2}                            & \tikzmark{3}              & \tikzmark{4}                            \\
		Early MArm.       &               & \tikzmark{5}berem                       & \tikzmark{6}berem         & \tikzmark{7}berem                       \\
		                  &               & \tikzmark{a}\cellcolor{gray!25}ku berem & \tikzmark{g}berem         & \tikzmark{e}\cellcolor{gray!25}ku berem \\
		                  &               &                                         &                           &                                         \\
		{[dialect split]} & \tikzmark{c}  & \tikzmark{b}                            & \tikzmark{h}              &                                         \\
		                  &               &                                         &                           &                                         \\
		Late MArm.        &               & \cellcolor{gray!25}berum em             & \tikzmark{k}berem         & \tikzmark{f}ku berem                    \\
		\textsc{mea}               &               & berum em                                & berem                     & kberem                                  \\\cline{3-5}
		                  & \tikzmark{i}  &                                         & \tikzmark{l}              &                                         \\
		\textsc{mwa}               & \tikzmark{c1} & \tikzmark{d}kə berem                    & \tikzmark{j}berem         & \cellcolor{gray!25}piti berem           \\
		\lspbottomrule    &               &                                         &                           &
	\end{tabularx}




\begin{tikzpicture}[overlay, remember picture, shorten >=.5pt, shorten <=.5pt, transform canvas={yshift=.25\baselineskip}]
	\draw [-] ([xshift=20pt,yshift=-7pt]{pic cs:a}) to ([xshift=20pt,yshift=-.75pt]{pic cs:b});
	\draw [-] ([xshift=20.75pt]{pic cs:b}) to ([xshift=0pt]{pic cs:c});
	\draw [-] ([yshift=.75pt]{pic cs:c}) to ([yshift=-.5pt]{pic cs:c1});
	\draw [->] ([xshift=0pt]{pic cs:c1}) to ([xshift=0pt]{pic cs:d});
	
	\draw [dotted,->] ([xshift=20pt,yshift=-7pt]{pic cs:e}) to ([xshift=20pt,yshift=7pt]{pic cs:f});
	
	\draw [dashed,->] ([xshift=12pt,yshift=-7pt]{pic cs:g}) to ([xshift=12pt,yshift=5pt]{pic cs:k});
	\draw [dashed] ([xshift=12.75pt,yshift=-3]{pic cs:h}) to ([xshift=3pt,yshift=-3pt]{pic cs:c});
	\draw [dashed] ([xshift=3pt,yshift=-3]{pic cs:c}) to ([xshift=3pt,yshift=0pt]{pic cs:i});
	\draw [dashed] ([xshift=3pt,yshift=0]{pic cs:i}) to ([xshift=12pt,yshift=0pt]{pic cs:l});
	\draw [dashed,->] ([xshift=12pt,yshift=0pt]{pic cs:l}) to ([xshift=12pt,yshift=7pt]{pic cs:j});

	
	\draw [-] ([xshift=12pt,yshift=-7pt]{pic cs:1}) to ([xshift=12pt,yshift=-.75pt]{pic cs:2});
	\draw [-] ([xshift=11.25pt,yshift=0pt]{pic cs:2}) to ([xshift=12.75pt,yshift=0pt]{pic cs:4});
	\draw [->] ([xshift=12pt,yshift=.75pt]{pic cs:2}) to ([xshift=12pt,yshift=5pt]{pic cs:5});
	\draw [->] ([xshift=12pt,yshift=.75pt]{pic cs:3}) to ([xshift=12pt,yshift=5pt]{pic cs:6});
	\draw [->] ([xshift=12pt,yshift=.75pt]{pic cs:4}) to ([xshift=12pt,yshift=5pt]{pic cs:7});
\end{tikzpicture}
\end{table}

As the table suggests, while both varieties have undergone very similar processes of retention and innovation, using different form-function pairs as their basis, it seems clear that, at least as morphosyntax is concerned, there was a significant varietal split in the Middle Armenian period when Eastern and Western varieties ceased to develop along the same lines; it is, however, at present not possible to identify a probable time period for this split owing to the insufficient availability of data and research already mentioned above.

\subsection{Hypoanalysis}
Quite apart from the particular diachrony of the present, subjunctive and future between Old and Modern Armenian, the developments outlined above are also of interest from a typological perspective in that they illustrate the interplay of ``regular'' grammaticalisation processes and other diachronic mechanisms of change. The concept of hypoanalysis in particular, as defined by \citep[126–127]{Croft2000}, is of interest in this context:

\begin{quote}
	In \textsc{hypoanalysis}, the listener reanalyzes a contextual semantic/functional property as an inherent property of the syntactic unit. In the reanalysis, the inherent property of the context […] is then attributed to the syntactic unit, and so the syntactic unit in question gains a new meaning or function.
\end{quote}

The keyword in this definition is the notion of contextuality. Where certain functions or semantics are originally caused by contextual clues such as adverbs, noun phrases or prepositional phrases,\footnote{In the present context, that is futurity, examples like \textit{tomorrow}, \textit{next week}, or \textit{in the coming days} come to mind.} hypoanalysed forms do not require these clues any more; in logical terms, A\&B$\rightarrow$C is simplified to A$\rightarrow$C, where A is the form in question, B the original context clue, and C the interpretation.

Hypoanalysis cannot, logically, apply in cases of complete innovation or creation \textit{ex nihilo}, as contextual ambiguities are not at play here. So, when early Middle Armenian creates the \textit{ku-}present, or late Middle Armenian the \textit{-um}-present, or Modern Western Armenian the \textit{piti}-future, these are all clear instances of grammaticalisation: original periphrases with specific semantics (progressive, deontic) are reanalysed, semantically bleached in the process, and retain only a functional value as well as being obligatorified for the relevant category (present, future). While other common processes associated with grammaticalisation are not as evidently at play here (phonetic reduction, decategorialisation), the concerned forms fit the ``grammaticalisation cline'' bill rather well, since progressives are known to turn into (imperfective) presents,\footnote{Cf., e.g., \citet{Deo2015}.} and deontic expressions into futures.

By contrast, a hypoanalytical process is at play in the functional specification of the old present as a subjunctive in early Middle Armenian. There, a set of forms sets off being functionally ambiguous (present, subjunctive, future) as a result of the loss of previous, more specific formations (the Old Armenian subjunctive and future periphrasis). As mentioned in §\ref{sec:meyer:emarm} above, this ambiguity could be contextually reduced by means of appropriate adverbs for the future, subjunctions, or periphrases for the subjunctive (e.g. \textit{t‘oł} + finite verb), whereas otherwise unmarked forms were interpreted as presents.

\largerpage[1.5]

As early Middle Armenian and later Modern Western Armenian clearly indicate, the \textit{ku-}periphrasis over time grammaticalised as the standard present tense. For a time, however, there must have been a meaningful difference between these periphrastic forms and the old presents; given the occasional future use of the periphrasis in early Middle Armenian and its later grammaticalisation in late Middle Armenian as a future tense, the particular semantics of \textit{ku} may have been neither purely indicative, as \citet{Karst1901} holds, nor originally progressive as per \citet{Vaux1995}.

For a time, then, the \textit{ku-}present and old present co-existed, both referring (in potentially aspectually different ways) to the present indicative. This contrast is eliminated with the grammaticalisation of the \textit{ku-}present as an aspectually unmarked form. The reduction of the old present to subjunctival (and future) function is, however, not simply a side-effect, since nothing in principle would have prevented the maintenance of an aspectual differentiation. Rather, 
a separate but concomitant hypoanalysis process must be assumed, chiefly because the contextually unmarked reading of the old present shifted to the subjunctive, which previously required contextual cues.

The cyclicity of grammaticalisation phases comes to the fore when a very similar process, the grammaticalisation of the \textit{-um-}present in late Middle Armenian, causes or coincides with the hypoanalysis of the \textit{ku-}present as a future. This development suggests that the \textit{ku-}form had retained a contextual future meaning in the first round of grammaticalisation (now decontextualised as its only function); likewise, this suggests that the old present gradually lost not only its present but also future meaning in the first round of hypoanalysis.

But why is it necessary – or at least preferable – to separate the two processes, that is the grammaticalisation of a periphrasis as a new tense and the specialisation of an older form by hypoanalysis, rather than seeing the latter as a side-effect of the former? Firstly, innovative forms can in principle lead to three fates for older homosemous forms: loss, maintenance, or specialisation. The Romance future, originally a deontic periphrasis consisting of an infinitive + \textit{have}, led to the loss of the established Latin \textit{b-}future; the same future forms as well as the English \textit{will-}future have been maintained for some time despite the coexistence and near-homosemy with periphrases of the \textit{going-to}-type; or the original forms can retain a specialised meaning, usually assuming pre-existing polysemy, as is the case in the Armenian example discussed here.

The notion of the side-effect could, of course, be maintained by adding polysemy as a conditioning factor. The data presented here are not enough to suggest that polysemy is a necessary condition for this side-effect specialisation, but taken together with the developments presented in §\ref{sec:meyer:modarm} below, it becomes clear that it isn't a sufficient condition. For in Modern Eastern Armenian, the \textit{k-}future is polysemous, but retains as its most elementary function the expression of an (aspectually unmarked) future, whereas a plethora of aspectually diverse future periphrases exist concurrently. This, then, is the second point: specialisation is not automatic, even in very similar circumstances.

The observations made here therefore complicate somewhat the explanations of the Armenian data as proposed by \citet{Haspelmath1998a} and \citet{BybeePerkinsPagliuca1994}, without however invalidating them. Instead, it is simply argued that the processes described there, that is the gain of futures and subjunctives from old presents, may not be automatic developments, conditioned or not, but separate, concomitant processes complementary to the grammaticalisation of innovative forms. In the case of the Armenian data, the simplest assumption would seem to be that this process was one of hypoanalysis, given the pre-existing polysemy of the old forms in both rounds.


	%Note that most morphosynt info is now contained in ku, as the prs remains ambiguous
	\section{Diversification in Modern Eastern Armenian}\label{sec:meyer:modarm}
After the slightly more theoretical consideration of the diachronic processes involved in the establishment of the Modern Armenian future and other basic tenses as described above, the situation in Modern Eastern Armenian in particular deserves closer attention.\footnote{In most contexts, including the present one, Modern Eastern Armenian is perhaps best understood as the more or less standardised, literary variety of the Eastern or  \textit{-um}-dialects. This literary variety began to take shape in the middle of the 19th century; Xač‘atur Abovyan's \textit{Verk‘ Hayastani} (‘Wounds of Armenia’, 1841) is traditionally considered the first piece of Modern Eastern Armenian literature.} There, next to the \textit{k(u)-}future already laid out in some detail, a set of other, more or less grammaticalised periphrases with future reference have established themselves. This section aims to briefly describe the form and function of these periphrases and to propose some first elements of a diachrony, without however entering into the finer details of the various constructions' semantics or complementary distribution.

These future periphrases occur in the context of a number of aspectual contrasts that pervade the Armenian verbal system, for which \citet[193–199]{DumTragut2009} recognises the following pairs:
\begin{itemize}
	\item perfect vs imperfective – e.g. perfective \textit{grec‘i} ‘I wrote’ vs imperfective \textit{grum ēi} ‘I was writing’;
	\item habitual-iterative vs processual – both expressed by the present or imperfect tense and differentiated adverbially;
	\item actional vs stative – the former being expressed by durative periphrases with the participle in \textit{-lis}.\footnote{This explanation of the Armenian TAME system is both incomplete and somewhat self-contradictory: the \textit{-lis-}participles, for instance, are processual rather than durative; the stative-resultative participle notion, as expressed by the \textit{-ac-}participle is not integrated; and expressions of past evidentiality are not clearly integrated \citep{Kozintseva2000}. For the present purpose, however, this does not pose an insurmountable problem.}
\end{itemize}
These aspectually marked future periphrases occur alongside the historic \textit{k-}future, which has been retained.\footnote{Modern grammatical writing in Armenian refers to these forms as the \textit{conditional mood} (Arm. \textit{paymanakan ełanak}) owing to its secondary, non-futuric uses. For the present purpose, the term \textit{k-}future will be used for ease of comprehension.} This synthetic tense, it would appear, is aspectually unmarked and next to its future use can also relate other notions such as generalised statements (\ref{ex:meyer:modarm1}) and potentiality (\ref{ex:meyer:modarm2}).

\ea \label{ex:meyer:modarm1}
\gll k‘ami=n k-p‘č‘i anjrew=ə k-teła bayc‘ leṙ=ə k-mna kangun.\\
wind.\textsc{nom=det} \textsc{fut-}blow.\textsc{3sg} rain.\textsc{nom=det} \textsc{fut-}fall.\textsc{3sg} \textsc{conj} mountain.\textsc{nom=det} \textsc{fut-}stay.\textsc{3sg} standing\\
\glt ‘The wind blows, the rain falls, but the mountain remains standing.’ \citep[254–255]{DumTragut2009}
\z \il{Modern Eastern Armenian}

\ea \label{ex:meyer:modarm2}
\gll ays hin patmut‘yun=ə deṙ k-lses gyuł-i tarec mardkanc‘ic‘\\
\textsc{dem} old story.\textsc{acc=det} still \textsc{fut-}hear.\textsc{2sg} village-\textsc{gen} aged man.\textsc{abl.pl}\\
\glt ‘This old story you may still hear from old village folk’ \citep[255]{DumTragut2009}
\z \il{Modern Eastern Armenian}

The periphrastic future expressions, too, can be constructed with forms of the \textit{k-}future, marked on the auxiliary verb \textit{linem} ‘to become’, and with the aid of participles (here called \textit{k-}periphrases). Parallel forms are, however, also available formed with the present tense of the auxiliary \textit{em} ‘to be’, the dynamic participle in \textit{-lu} and other participles (here called  \textit{-lu-}periphrases). \tabref{tab:meyer:modarmfut} gives an overview of these forms.\footnote{The terms \textit{dynamic}, \textit{stative}, and \textit{processual} used here reflect the terminology in \citet{DumTragut2009}.}

\begin{table}
	
	\caption{Overview of future expressions in Modern Eastern Armenian}\label{tab:meyer:modarmfut}
	\begin{tabularx}{0.8\textwidth}{Xll}
		\lsptoprule    & \textit{-lu}-periphrases   & \textit{k-}periphrases                         \\ \midrule
		dynamic        & \textit{berelu em}        & \textit{berelu klinem}  \\
		               & ‘I shall carry’           &        ‘I shall be about to carry’                                            \\
		stative        & \textit{paṙkac em linelu} & \textit{paṙkac klinem}      \\
		               & ‘I shall be lying down’   &        ‘I shall be lying down’                                            \\
		processual     & \textit{gnalis em linelu} & \textit{gnalis klinem}           \\
		               & ‘I shall be going’        &          ‘I shall be going’                                          \\
		\lspbottomrule 
	\end{tabularx}
\end{table}

As the terminology and glosses suggest already, these periphrases refer to immediate and/or planned actions (dynamic), resulting states (stative),\footnote{These forms are usually reserved for position verbs.} and continuous actions (processual). Examples (\ref{ex:meyer:modarm3}–\ref{ex:meyer:modarm7}) illustrate these different uses briefly.


\ea \label{ex:meyer:modarm3}
\gll Lewon, varə … meładrakan t‘łt‘i=d meǰ hastat im anun=ə grelu es, č‘-ē?\\
	\textsc{pn} tomorrow {} accusatory letter.\textsc{dat.sg-2poss} in surely \textsc{1sg.poss} name.\textsc{acc.sg=det} write.\textsc{ptcp.dyn} be.\textsc{2sg.prs}, \textsc{neg-}be.\textsc{3sg.prs}\\
	\glt ‘Lewon, tomorrow […] you are surely going to mention my name in your accusatory letter, aren't you?’ (\textit{Lur č‘ka},  Gurgen Xanǰyan, 2016) – dynamic \textit{-lu-}periphrasis
	\z \il{Modern Eastern Armenian}

\ea \label{ex:meyer:modarm4}
\gll erewum ē du stipvac es linelu aysteł arazel {mi k‘ani žam}, et‘e oč‘ ambołǰ gišer=ə\\
	seem.\textsc{prs.impf} be.\textsc{3sg.prs} \textsc{2sg.nom} force.\textsc{pass.ptcp.res} be.\textsc{2sg.prs} become.\textsc{ptcp.dyn} here graze.\textsc{inf} for-some-time if \textsc{neg} all evening=\textsc{det}\\
	\glt ‘It seems you'll be forced to graze here for some time, if not the whole evening.’ (tr. of \textit{The Headless Horseman}, Mayne Reed, 1958) – stative \textit{-lu-}periphrasis
	\z \il{Modern Eastern Armenian}
	
\ea \label{ex:meyer:modarm5}
\gll Sranic‘ heto erb Erewanum arjanner teładrelu k-linenk‘, k-haskanank‘, or da šat patasxanatu u nurb gorc ē\\
\textsc{dem.abl.sg} after when \textsc{pn.loc.sg} statue.\textsc{acc.pl} erect.\textsc{ptcp.dyn} \textsc{fut-}become.\textsc{1pl} \textsc{fut-}understand.\textsc{1pl} \textsc{comp} \textsc{dem.nom.sg} very responsible and subtle work.\textsc{nom.sg} be.\textsc{3sg.prs}\\
\glt ‘After this, when we will be set to erect statues in Erewan, we will understand that this is a very responsible and subtle task’ (\textit{Aṙawōt}, 01.10.2002) – dynamic \textit{k-}periphrasis
\z \il {Modern Eastern Armenian}

\ea \label{ex:meyer:modarm6}
\gll Aynpes or, šat ban kaxvac k-lini ayd k‘vearkut‘yan ardyunk‘ic‘\\
thus \textsc{comp} many think.\textsc{nom.sg} depend.\textsc{ptcp.res} \textsc{fut-}become.\textsc{3sg} \textsc{dem} vote.\textsc{gen.sg} outcome.\textsc{abl.sg}\\
\glt ‘Thus a lot will depend on the outcome of that vote.’ (\textit{Aṙawōt}, 01.06.2004) – stative \textit{k-}periphrasis
\z \il{Modern Eastern Armenian}

\ea \label{ex:meyer:modarm7}
\gll erb kardalis \textsc{k-}lines, ušadrut‘yun darju Trimalk‘ioni mecark‘i vra\\
when read.\textsc{ptcp.proc} become.\textsc{2sg} attention.\textsc{acc.sg} turn.\textsc{2sg.imv} \textsc{pn.gen.sg} banquet.\textsc{gen.sg} upon\\
\glt ‘When you will be reading, turn your attention to the \textit{Cena Trimalchionis}’ (Henrik Senkewič‘, \textit{Yo Ert‘as}, 1983) – processual \textit{k-}periphrasis
\z \il {Modern Eastern Armenian}
	
This aspectual richness of expression, surprising in view of the paucity and simplicity of future formations in later Middle Armenian, deserves to be studied in its own right as regards its use and distribution. For the present purpose, that is an overview of the diachrony of Armenian futures, a glimpse at the development of these periphrases over time will have to suffice. With the help of the \textit{Eastern Armenian National Corpus} (EANC),\footnote{The corpus, accessible at \href{http://eanc.net}{http://eanc.net}, contains more than 110,000,000 tokens across written discourse (prose, poetry; fictional, non-fictional), oral discourse and press archives and includes content from the middle of the 19th century onwards.} the figures of absolute occurrences in the whole corpus of these periphrases and their first attestations have been gathered in \tabref{tab:meyer:modarmfut2}.

\begin{table}
	
	\caption{Occurrences and first attestations of future periphrases in Modern Eastern Armenian\newline The figures in parentheses indicate the number of texts in which the periphrasis occurs; if two years are cited, there is only one occurrence attested for the first year.}\label{tab:meyer:modarmfut2}
	\begin{tabularx}{1\textwidth}{Xrrrr}
		\lsptoprule
		Type         & \textit{-lu-}periphrases & Year & \textit{k-}periphrases  & Year      \\ \midrule
		dynamic (-\textit{lu})     & 71,300 (7,400)                   & 1860 & 40 (39)                        & 1870/1955 \\
		stative (-\textit{ac})     & 1,289 (1,033)                    & 1958 & 8,662 (3,262)                  & 1841      \\
		processual (-\textit{lis}) & 2 (2)                            & 2005 & 103 (66)                       & 1870/1950 \\
		\lspbottomrule          
	\end{tabularx}
\end{table}

On a purely numerical basis, the data indicate that three of the periphrases (dynamic \textit{k-}, processual \textit{-lu-}, and processual \textit{k-}) are marginal occurrences and comparatively recent creations. The reason for the low occurrence of processual futures is, perhaps, self-explanatory: they are aspectually very precise but describe an action that has not yet happened and, therefore, cannot be predicted – or need not be specified to such a degree of precision. The same is true, \textit{mutatis mutandis}, for the dynamic \textit{k-}periphrasis which describes a planned or imminent action in the future and is similarly aspectually over-specific for common use.

That leaves three more frequent periphrases. The most common one – the dynamic \textit{-lu-}pe\-ri\-phrasis (\textit{berelu em}) – is also the morphologically simplest and in all likelihood the best analogy to the pairing of synthetic futures and \textit{going-to-}periphrases found in Romance and Germanic. Alongside it, the stative \textit{k-}periphrasis is attested similarly early, even though only about 11\% as often; this may be related to its aspectually more specific function and the smaller number of verbs with which it is appropriate. Finally, the stative \textit{-lu-}periphrasis is a more recent innovation and, perhaps for that reason, again not attested as widely as the corresponding \textit{k-}periphrasis. All of these figures contrast with the quasi-ubiquitous simple \textit{k-}future (\textit{kberem}), which is attested 630'680 times over 9'290 texts.\footnote{These figures are, to a certain extent, a simplification, since the \textit{k-}periphrases can (but do not always) have additional modal connotations which have not been teased out here; cf. \citet{Kozintseva2005}. This distinction is, however, likely not relevant for the diachrony discussed here.}

In diachronic terms, then, and taking into account the data above, it seems most plausible to assume that next to the simple \textit{k-}future, the stative \textit{k-}periphrasis developed early to allow for the expression of the aspectual contrast of action vs state/result; at a similar time, the dynamic \textit{-lu-}periphrasis arose, perhaps to distinguish a different kinds of futurity (planned vs definite; immediate vs remote; etc.). The other, less frequent and later developed periphrases, are analogical extensions: the stative \textit{-lu-}periphrasis creates a parallel form to the stative \textit{k-}periphrases; the processual periphrases extend present-tense formations (type \textit{gnalis em}) to the future context. The fact that the processual \textit{-lu-}periphrasis in particular occurs exclusively in grammar books raises the question to what extent all three of these less-common periphrases are the result of paradigmatisation without grammaticalisation. 

It goes without saying that the present sketch of the future expressions in Modern Eastern Armenian does not do justice to its complexity, both in synchronic and in diachronic terms. It does serve, however, to give a brief overview of these formations and to raise questions for the future, in particular the following: Is the development of the aspectual specificity –  hinted at in the future periphrases and more developed in other tenses – a purely internal development or, at least in part, informed by or modelled on similar distinctions in contact languages?

	\section{Conclusions and perspectives}\label{sec:meyer:conc}
	The aim of this paper was to describe the diachrony of future expressions in Armenian, from its oldest attestations to the present varieties of the language. Beginning with a polysemous subjunctive, the Armenian language innovated a separate future category only in its Medieval stages, after the broader dialect split into an Eastern and Western variety had taken place. The Western varieties separately grammaticalised an originally deontic periphrasis as their future tense, whereas the Eastern varieties repurposed a present periphrasis. In Modern Eastern Armenian, this same future tense is maintained; at the same time, more aspectually specific periphrases have been created – some more grammaticalised than others — to analogously extend to the future the aspectual specificity Armenian exhibits in other tenses.
	
	Over the course of its development, future expressions have had regular and close interactions both with the present tense and the subjunctive. In particular, it has been argued that on two occasions, the Armenian tense/mood categories have undergone a pair of related but separate grammaticalisation and hypoanalysis processes – first for the Old Armenian present turning into a subjunctive; then the late Middle Armenian present turning into a future. Different processes, namely analogical extension and paradigmatisation, are at play in the development of the future periphrases in Modern Eastern Armenian.
	
	Apart from illustrating the development of the future in Armenian and thereby providing a detailed example of a fairly complex set of innovations, this paper has sought to complicate to a certain extent the well-known accounts of developments in Armenian cited in the linguistic literature as well as the relationship between hypoanalysis and grammaticalisation. Yet, as has been pointed out, much work remains to be done, in particular as regards the Middle Armenian data, to gain a complete understanding of the changes happening there and their potential interactions. Likewise, the system of future expressions in Modern Armenian deserves more attention and should prove a fruitful ground for research into grammaticalisation and related processes as well as into the diachrony of the Armenian aspect system.
	
	
	
	
	
	
	
	

	
	\section*{Abbreviations}
	\begin{tabularx}{.5\textwidth}{@{}lQ@{}}
		\textsc{1} & first person\\
		\textsc{2} & second person\\
		\textsc{3} & third person\\
		\textsc{abl} & ablative\\
		\textsc{acc} & accusative\\
		\textsc{aor} & aorist \\
		\textsc{comp} & complementiser\\
		\textsc{conj} & conjunction\\
		\textsc{dat} & dative\\
		\textsc{dem} & demonstrative pronoun\\
		\textsc{det} & determiner\\
		\textsc{dyn} & dynamic\\
		\textsc{emph} & emphatic\\
		\textsc{f} & feminine\\
		\textsc{fut} & future\\
		\textsc{gen} & genitive\\
		\textsc{imv} & imperative\\
		\textsc{indf} & indefinite\\
		\textsc{inf} & infinitive\\
		\textsc{ins} & instrumental\\
		\textsc{loc} & locative\\
		
	\end{tabularx}%
	\begin{tabularx}{.5\textwidth}{@{}lQ@{}}
		\textsc{m} & masculine\\
		\textsc{\textsc{mea}} & Modern Eastern Armenian\\
		\textsc{\textsc{mwa}} & Modern Western Armenian\\
		\textsc{nom} & nominative \\
		\textsc{obj} & object(-marker)\\
		\textsc{perl} & perlative\\
		\textsc{pl} & plural \\
		\textsc{pn} & proper noun\\
		\textsc{poss} & possessive\\
		\textsc{proc} & processual\\
		\textsc{prog} & progressive\\
		\textsc{prs} & present\\
		\textsc{pst} & past\\
		\textsc{ptc} & particle\\
		\textsc{ptcp} & participle\\
		\textsc{refl} & reflexive\\
		\textsc{rel} & relative pronoun\\
		\textsc{res} & resultative\\
		\textsc{sbjv} & subjunctive\\
		\textsc{sg} & singular\\
		\textsc{vbadj} & verbal adjective\\
	\end{tabularx}
	
	\section*{Acknowledgements}
	I'm very grateful to the editors, Stefan Hartmann and Lena Schnee, for the organisation of the conference that led to the publication of this volume, and for all participants for the interesting discussions of this subject. Further thanks go to Neige Rochant for her thought-provoking questions and thoughts on this topic at a departmental seminar, and to the helpful and constructive comments of the eagle-eyed peer-reviewers. All errors and omissions are, of course, mine.
	
	\sloppy
	\printbibliography[heading=subbibliography,notkeyword=this]
\end{document}
